\documentclass{article}
\usepackage{amsmath, amssymb, amsthm}

\setlength{\parindent}{0pt}

\newtheorem*{theorem}{Theorem}

\begin{document}

\begin{theorem}
    The second player has a winning strategy if and only if the number of sticks
    equals $4k+1$ for some $k \in \mathbb{N}$; otherwise, the first player wins.
\end{theorem}

\begin{proof}
    The proof is by strong induction. Let $P(n)$ denote the claim for $n$ sticks.

    \medskip
    \textbf{Base cases.}
    \begin{description}
        \item[$n=1$:] $P_2$ wins immediately.
        \item[$n=2$:] $P_1$ takes 1; $P_1$ wins.
        \item[$n=3$:] $P_1$ takes 2; $P_1$ wins.
        \item[$n=4$:] $P_1$ takes 3; $P_1$ wins.
        \item[$n=5$:] $P_1$ must leave 2, 3, or 4 sticks, so $P_2$ wins by the above.
    \end{description}

    \medskip
    \textbf{Inductive step.}
    Assume $P(1), P(2), \ldots, P(n)$ hold for some $n \geq 5$; we show $P(n+1)$.

    \medskip
    \textit{Case 1: $n+1 = 4k+1$.}
    After $P_1$'s move, the pile has $4k$, $4k-1$, or $4k-2$ sticks.
    In each subcase $P_2$ can restore the pile to the form $4(k-1)+1$:
    \begin{description}
        \item[$4k$ sticks:] $P_2$ takes 3, leaving $4(k-1)+1$.
        \item[$4k-1$ sticks:] $P_2$ takes 2, leaving $4(k-1)+1$.
        \item[$4k-2$ sticks:] $P_2$ takes 1, leaving $4(k-1)+1$.
    \end{description}
    By the inductive hypothesis $P_2$ wins from $4(k-1)+1$, so $P_2$ wins from $4k+1$.

    \medskip
    \textit{Case 2: $n+1 \neq 4k+1$.}
    Then $n+1$ is one of $4k+2$, $4k+3$, or $4k+4$.
    In each subcase $P_1$ can move directly to $4k+1$:
    \begin{description}
        \item[$4k+2$ sticks:] $P_1$ takes 1, leaving $4k+1$.
        \item[$4k+3$ sticks:] $P_1$ takes 2, leaving $4k+1$.
        \item[$4k+4$ sticks:] $P_1$ takes 3, leaving $4k+1$.
    \end{description}
    By the inductive hypothesis $P_2$ wins from $4k+1$, so $P_1$ wins from each of these.

    \medskip
    In both cases $P(n+1)$ holds, completing the induction.
\end{proof}

\end{document}
